\chapter{Cái đúng có giá}
\markboth{Cái đúng có giá}{Right Has a Cost}

\section{Một học sinh không quay cóp nên điểm thấp hơn bạn.}
\begin{truyen}{Một học sinh không quay cóp nên điểm thấp hơn bạn.}{Doing right can look like losing first.}
\chuhoa{E}{m biết mình có thể nhìn sang bài bạn mà không ai phát hiện. Nhưng em không làm. Kết quả thấp hơn làm em buồn, nhưng tối đó em ngủ yên. Có những cái đúng khiến ta thua trước mắt nhưng đỡ phải khinh mình về sau.}
\textit{(The student knew cheating was possible without getting caught. Yet he did not do it. The lower result hurt, but that night he slept peacefully. Some right choices make us lose in the short term and save us from despising ourselves later.)}
\end{truyen}
\begin{baihoc}
Cái đúng nhiều khi không thưởng ngay, nhưng nó giữ phần cốt lõi của con người.
\textit{(What is right often does not reward immediately, but it protects the core of a person.)}
\end{baihoc}
\begin{ghinhoanh}\textbf{Short English to remember:}

Doing right can look like losing first.
\end{ghinhoanh}
\begin{tuhocnhanh}
1. Em có đang chịu một thiệt thòi nào vì không làm sai không? \textit{(Are you enduring any disadvantage because you refuse to do wrong?)}

2. Điều đó có làm em vững hơn về sau không? \textit{(Will that make you stronger later?)}
\end{tuhocnhanh}
\ngancach

\section{Một nhân viên nói thật nên mất lòng sếp.}
\begin{truyen}{Một nhân viên nói thật nên mất lòng sếp.}{Truth can be costly.}
\chuhoa{A}{nh biết bản báo cáo có vấn đề và không muốn ký xác nhận qua loa. Việc đó khiến cấp trên khó chịu. Anh không thoải mái ngay lúc ấy, nhưng anh biết cái giá của sự thật vẫn rẻ hơn cái giá của việc đồng lõa.}
\textit{(He knew the report had problems and refused to sign it carelessly. That upset his superior. It was uncomfortable in the moment, but he knew the price of truth was still cheaper than the price of complicity.)}
\end{truyen}
\begin{baihoc}
Nói thật có thể làm khó ta một lúc; nói sai theo đám đông có thể làm hại ta rất lâu.
\textit{(Speaking truth may trouble us for a while; following a false crowd may damage us for a long time.)}
\end{baihoc}
\begin{ghinhoanh}\textbf{Short English to remember:}

Truth can be costly, but lies cost deeper.
\end{ghinhoanh}
\begin{tuhocnhanh}
1. Em có đang im lặng trước điều mình biết là không đúng không? \textit{(Are you staying silent about something you know is wrong?)}

2. Em sợ mất điều gì nếu nói thật? \textit{(What are you afraid of losing if you speak honestly?)}
\end{tuhocnhanh}
\ngancach

\section{Một người không chịu lách luật nên chậm giàu hơn.}
\begin{truyen}{Một người không chịu lách luật nên chậm giàu hơn.}{Clean money often grows slower.}
\chuhoa{A}{nh nhìn người khác đi nhanh hơn và đôi lúc thấy dao động. Nhưng mỗi lần nghĩ tới hậu quả của tiền sai cách, anh lại chấp nhận đi chậm. Đúng đường không phải lúc nào cũng nhanh, nhưng nó ít để lại nỗi sợ hơn.}
\textit{(He watched others move faster and sometimes felt shaken. Yet each time he considered the cost of dirty money, he accepted a slower pace. The right road is not always fast, but it leaves less fear behind.)}
\end{truyen}
\begin{baihoc}
Đi chậm vì giữ sạch mình còn hơn đi nhanh bằng thứ một ngày nào đó sẽ quay lại đòi giá.
\textit{(Moving slowly while staying clean is better than moving fast with something that will later demand payment.)}
\end{baihoc}
\begin{ghinhoanh}\textbf{Short English to remember:}

Clean money often grows slower.
\end{ghinhoanh}
\begin{tuhocnhanh}
1. Em có đang sốt ruột vì người khác đi nhanh hơn không? \textit{(Are you becoming impatient because others move faster?)}

2. Em muốn giữ điều gì sạch nhất trong đời mình? \textit{(What do you most want to keep clean in your life?)}
\end{tuhocnhanh}
\ngancach

\section{Một cô gái dám từ chối điều trái lòng dù bị chê cứng.}
\begin{truyen}{Một cô gái dám từ chối điều trái lòng dù bị chê cứng.}{Boundaries can be right.}
\chuhoa{K}{hông phải lời từ chối nào cũng là lạnh lùng. Có lần cô nói không vì biết nếu chiều theo sẽ tự làm đau mình. Người khác không thích, nhưng chính sự rõ ràng đó giữ cô khỏi một chuyện sai với chính mình.}
\textit{(Not every refusal is cold. Once she said no because she knew agreeing would harm herself. Others did not like it, but that clarity kept her from betraying herself.)}
\end{truyen}
\begin{baihoc}
Sống đúng không chỉ là tử tế với người khác, mà còn là không phản bội giới hạn lành mạnh của mình.
\textit{(Living rightly is not only being kind to others, but also not betraying your own healthy boundaries.)}
\end{baihoc}
\begin{ghinhoanh}\textbf{Short English to remember:}

Boundaries can be right.
\end{ghinhoanh}
\begin{tuhocnhanh}
1. Em có đang đồng ý quá nhiều vì sợ bị chê không? \textit{(Are you saying yes too often because you fear criticism?)}

2. Ranh giới đúng nào em cần giữ chắc hơn? \textit{(What right boundary do you need to hold more firmly?)}
\end{tuhocnhanh}
\ngancach

\section{Một người dám nhận lỗi trước đám đông.}
\begin{truyen}{Một người dám nhận lỗi trước đám đông.}{Admitting wrong protects the right thing.}
\chuhoa{K}{hi sai sót xảy ra, anh có thể đổ qua hoàn cảnh. Nhưng anh nhận phần trách nhiệm của mình trước mọi người. Khoảnh khắc ấy làm anh thấy xấu hổ, nhưng cũng giữ cho tập thể còn tin vào sự ngay thẳng.}
\textit{(When a mistake happened, he could have blamed circumstances. Instead, he accepted his part publicly. The moment was embarrassing, but it preserved the group's trust in honesty.)}
\end{truyen}
\begin{baihoc}
Nhận mình sai không làm cái đúng mất đi; nó đưa ta quay lại gần cái đúng hơn.
\textit{(Admitting you are wrong does not destroy what is right; it brings you closer to it.)}
\end{baihoc}
\begin{ghinhoanh}\textbf{Short English to remember:}

Admitting wrong protects what is right.
\end{ghinhoanh}
\begin{tuhocnhanh}
1. Em có đang chậm nhận lỗi vì ngại xấu hổ không? \textit{(Are you delaying an apology because of embarrassment?)}

2. Nếu nhận ngay, điều gì sẽ được giữ lại? \textit{(If you admitted it now, what could still be preserved?)}
\end{tuhocnhanh}
\ngancach